\section{Weak product testing}\label{sec:product-tests}

The preceding histogram lemma gives compact families within each
leaf. We now prove a different statement: a fixed global test on a
tuple of keys can be replaced, for the integrations that occur here,
by a bounded sum of products of global single-key functions. Both its
reference and its actual-query conclusions are needed.

Fix a query template with $r$ positions and with labels distinct at
each used endpoint. Its input laws and reference measures are those
of \Cref{sec:histograms}. In particular
$\overline\Omega_n=\bigotimes_{j=1}^r\Omega_{l_j,q_j,n}$ and
$\nu_n$ is its off-position-Gram-conditioned reference on $\cK_n$.
The additional prescribed binary tests are the within-endpoint
ordered $L,R$ evaluations from \Cref{sec:realization}. Let
$\mathcal A_n(O,z)$ be the successful binary indicator from
\Cref{lem:binary-mixing}, including the off-position Gram zeros.
All common functions below inspect primal keys only. Channel maps
may enter the unary input weights and the orientation filters.

\begin{theorem}[Weak product comparison]\label{thm:weak-product}
Let $W_n$ be any sequence of real functions on $\cK_n$ with
$\norm{W_n}_\infty\le B$, where $B$ is fixed. Extend it to the product
of the single-key spaces with the same bound. For every $\eps>0$
there are constants $L,C<\infty$, independent of $n$, and functions
\begin{equation}\label{eq:weak-product-approximant}
 S_n(k_1,\ldots,k_r)
   =\sum_{a=1}^{L}c_{a,n}\prod_{j=1}^{r}f_{a,j,n}(k_j),
 \qquad
 \abs{c_{a,n}}\le C,\quad \abs{f_{a,j,n}}\le1,
\end{equation}
with the following properties.
\begin{enumerate}[label=(\roman*)]
\item For arbitrary global unary functions $g_{j,n}$ of bound one,
\begin{equation}\label{eq:weak-product-reference}
 \abs{\int(W_n-S_n)\prod_jg_{j,n}(k_j)\,d\nu_n}\le\eps
\end{equation}
for all sufficiently large $n$.
\item Under any mixed unit law satisfying
\eqref{eq:histogram-mixed-caps}, outside an exponentially small
exceptional set of orientations,
\begin{equation}\label{eq:weak-product-actual}
 \sup_{\abs{w_j}\le1}
 \abs{\E_z\mathcal A_n(O,z)
       (W_n-S_n)(K(O,z))\prod_jw_j(O,z_j)}\le\eps.
\end{equation}
The unary input weights may depend on the entire fixed orientation
and on any additional marks fixed with it. The exceptional set is
independent of these weights. Thus one may prescribe channel flags
such as $p_i$ in a unary weight without making them common key tests
or asserting positivity for their individual values.
\end{enumerate}
The functions $f_{a,j,n}$ are global key functions; they can be chosen
to be indicators of sets, or finite linear combinations of such
indicators with uniformly bounded complexity. They do not depend on
the sampled orientation. The constants may depend on $B,\eps$ and
the fixed template and construction parameters, but not on $n$.

An additional orientation-only filter of bound one may multiply
\eqref{eq:weak-product-actual} without changing it. In particular,
after averaging normalized leaves, the same error control applies to
unnormalized admissibility-filtered query measures, even when the
filter is chosen using the opposite leaf. No assertion is made for
an arbitrary additional joint input filter.
For distinct marked laws, use their fixed-weight mixture, with caps
checked after forgetting the marks. A retained ordered type pair or
different-colour condition is an orientation/leaf-pair filter, so the
same conclusion applies before normalization. All leaf averages use
their actual weights; no marginal hypothesis on a separately
normalized colour class is needed.
\end{theorem}

The proof uses a finite family of twists. If $I,J$ partition a set of
positions into two nonempty groups, let $\mathcal T_{I,J}$ consist of
all products of their available inter-position global Gram signs:
the dot products of a plus slot at one position with a minus slot at
another position in the same component. These signs are defined on
the product reference whether or not that Gram is prescribed by a
raw same-endpoint frame law. Every such twist equals one on
$\cK_n$. The number of twists is a constant depending only on the
query shape.

For a real function $D$ on the keys of $I\cup J$, define its twisted
cut norm by
\begin{equation}\label{eq:twisted-cut-norm}
 \norm{D}_{\square,\mathcal T}
 =\max_{\tau\in\mathcal T_{I,J}}
   \sup_{\abs{f},\abs{g}\le1}
   \abs{\int D(k_I,k_J)\tau(k_I,k_J)f(k_I)g(k_J)
                  \,d\overline\Omega_{I\cup J,n}}.
\end{equation}
Here $f,g$ may be arbitrary functions on their whole respective
groups. Only after the recursion below will they become unary.

We use the finite-partition energy-increment method underlying weak
regularity; see \citet{frieze1999}. The additional Gram conditioning
needed here is retained explicitly in the following self-contained proof.

\begin{lemma}[Finite twisted regularity]
\label{lem:product-twisted-regularity}
Let $\abs{W}\le B$ on the product key space of $I\cup J$, and let
$\delta>0$. There is a decomposition
\begin{equation}\label{eq:twisted-regularity-decomposition}
 W=\sum_{a=1}^{L'}c_a f_a(k_I)g_a(k_J)\tau_a(k_I,k_J)+D,
\end{equation}
where $\abs{f_a},\abs{g_a}\le1$, $\tau_a\in\mathcal T_{I,J}$,
$\abs{c_a}\le B$, $\norm{D}_\infty\le2B$, and
$\norm{D}_{\square,\mathcal T}\le\delta$. The number $L'$ is bounded
in terms of $B,\delta$ and the twist count only. The functions $f_a,g_a$
may be chosen as indicators of cells in finite group partitions.
\end{lemma}

\begin{proof}
Let $\Gamma$ be the finite list of crossing Gram bits generating the
twists. Begin with the trivial partitions of the two groups. For any
current pair of finite partitions, condition $W$ on their cells and
on $\Gamma$; call the resulting conditional expectation $S$. Then
$\norm{S}_\infty\le B$ and $\norm{W-S}_\infty\le2B$.

If $D=W-S$ has twisted cut norm greater than $\delta$, choose a
twist and witnesses $f,g$ giving correlation greater than $\delta$.
Quantize $f,g$ uniformly within their bounds finely enough, with a
mesh depending only on $B,\delta$, that their quantized product times
the twist still correlates with $D$ by more than $\delta/2$. Refine
the group partitions by the values of these quantized functions.
The witness product is measurable for the new conditioning algebra
and has $L^2$ norm at most one. Thus the new conditional expectation
of $D$ has $L^2$ norm at least $\delta/2$. Since $D$ is orthogonal
to the old conditioning algebra, the squared $L^2$ norm of the new
conditional expectation of $W$ increases by at least $\delta^2/4$.
It is bounded above by $B^2$. At most $4B^2/\delta^2+1$
refinements are possible. Each increases the partition sizes by a
factor bounded in terms of $B,\delta$, so the final partitions have
bounded size.

On each product of group cells, the final conditional expectation is
a bounded function of $\Gamma$. Assign value zero to any impossible
combination of a cell pair and Gram bits. Its finite Fourier
expansion on the Gram bits has coefficients of absolute value at
most $B$. Multiplying by the two cell indicators gives
\eqref{eq:twisted-regularity-decomposition}. The decomposition is
exact everywhere on the product key space, regardless of which Gram
bit values occur in a particular cell pair.
\end{proof}

We next prove that a small twisted-cut remainder is negligible on
actual queries. This is the step that is not supplied by unary
histograms alone.

\begin{lemma}[A terminal remainder on actual inputs]
\label{lem:product-terminal-remainder}
Let $I,J$ be a nontrivial partition of some of the template positions.
Let $D_n$ be a common function of their global keys with
$\norm{D_n}_\infty\le B_1$ and
$\norm{D_n}_{\square,\mathcal T}\le\delta$. Let $\psi_n(I,J)$
be any product of the specified nominal $L,R$ signs and global Gram
signs crossing from $I$ to $J$, evaluated on their actual inputs and
keys. For each fixed $\varepsilon_{\mathrm{emp}}>0$, except on exponentially small mixed
mass,
\begin{equation}\label{eq:product-remainder-actual}
 \sup_{\abs{f},\abs{g}\le1}
 \abs{\E_{z_I,z_J}D_n(K_I,K_J)\psi_n(I,J)f(O,z_I)g(O,z_J)}
 \le\bigl(C_3\delta+\varepsilon_{\mathrm{emp}}+o(1)\bigr)^{1/4}.
\end{equation}
The constant $C_3$ and the vanishing error depend only on $B_1$ and
the fixed template. The weights may depend arbitrarily on the fixed
orientation. The exceptional set does not depend on those weights.
\end{lemma}

\begin{proof}
For one fixed orientation, put
$U(I,J)=D_n(K_I,K_J)\psi_n(I,J)$ and use independent copies of the
parameter groups $I,J$. Cauchy--Schwarz in $I$ gives
\[
 \abs{\E U(I,J)f(I)g(J)}^2
 \le\E_I\bigl(\E_JU(I,J)g(J)\bigr)^2.
\]
Expand the square, and apply Cauchy--Schwarz in the two $J$ variables.
The factors $g(J_0)g(J_1)$ have bound one. We obtain
\begin{equation}\label{eq:product-four-corner}
 \abs{\E U(I,J)f(I)g(J)}^4
 \le Q_n(O):=
 \E_{I_0,I_1,J_0,J_1}
       \prod_{a,b\in\{0,1\}}U(I_a,J_b).
\end{equation}
In particular $Q_n(O)\ge0$: it is the expectation over $J_0,J_1$ of
the square of $\E_IU(I,J_0)U(I,J_1)$. All separate weights, including
their orientation dependence, have disappeared.

The right-hand side is a bounded common input-and-key test, involving
two fresh copies of each original position. It depends only on the
primal maps: the channel-dependent weights and all marks have
disappeared in \eqref{eq:product-four-corner}, and the nominal $L,R$
signs involve only the displayed fresh parameters. Repeated labels across
these copies are permitted by \Cref{lem:typicality}. It follows that,
outside an exponentially small mixed exceptional set,
\begin{equation}\label{eq:product-four-corner-typical}
 Q_n(O)\le\E_{\rawlaw^2}Q_n+\varepsilon_{\mathrm{emp}}.
\end{equation}
We bound the raw mean by $C_3\delta+o(1)$.

Fix the nominal inputs of the four groups. Their columns are jointly
independent within each endpoint, component, and mode except with
$2^{-\Omega(n)}$ probability, since all copies have fresh $Z$ inputs.
The dependent-input contribution is $o(1)$ by the fixed sup bound.
On the independence event, \Cref{lem:frame-images} gives the uniform
image law with prescribed same-endpoint Grams and injectivity.
Ignore injectivity at an exponentially small error and express this
law relative to the product of all the individual single-key
references. There are only boundedly many additional Gram bits, and
their normalization is bounded uniformly in the nominal inputs by
\Cref{lem:histogram-gram-normalization}. Expand their indicators into
characters. The nominal $L,R$ phases are now constants. The global
Gram phases already in the four copies of $\psi_n$ can be combined
with the expansion characters.

For completeness, distinguish a slot by its group $I$ or $J$, its
copy, its position, component, and sign. Remove the diagonals already
in the reference and collapse repeated descriptions of the same
unordered slot dot. Every remaining Gram coefficient belongs to
exactly one of the following classes:
\begin{enumerate}[label=(\alph*)]
\item both slots lie in one of $I_0,I_1,J_0,J_1$;
\item the slots join $I_0$ to $I_1$, or $J_0$ to $J_1$;
\item the slots lie in one of the four $I_a$--$J_b$ quadrants.
\end{enumerate}
Class (a) contributes bounded separate group phases. Suppose a
nonzero class-(b) pattern between $I_0,I_1$ remains. Fix both $J$
groups. The four $D_n$ factors split into two bounded functions of
$I_0$ and $I_1$. The class-(c) phases become separate functions of
those two groups, as do all class-(a) phases. Their only remaining
interaction is the nonzero coefficient rectangle between $I_0$ and
$I_1$, whose bit rank is at least $N$. The individual diagonal
densities may also be absorbed into the respective bounded weights.
\Cref{lem:walsh} bounds this term by $O(2^{-N/2})$. The same argument
applies to a nonzero $J_0$--$J_1$ pattern.

If neither class-(b) pattern remains, attach each class-(c) pattern
as an available twist to its corresponding $D_n(K_{I_a},K_{J_b})$.
Fix $I_1,J_1$. The remaining integral in $I_0,J_0$ has one factor
$D_n$ times an available twist, tested against bounded separate
functions. The two other unfixed $D_n$ factors have bound $B_1$ and
the fixed factor has bound $B_1$. The twisted cut estimate therefore
bounds this term by $B_1^3\delta$, with the harmless fixed
normalization factors included in the constant.

These cases exhaust the character expansion. The groups $I$ and $J$
may each contain positions from both endpoints of the unit: this only removes some
same-endpoint Gram bits from the expansion and does not create a new
class. A nonzero coefficient rectangle has rank at least one in
distinguished global slots, hence bit rank at least $N$.
The two dots $P_a\cdot Q_b$ and $P_b\cdot Q_a$, for example, use
different slots and are not duplicates. The prescribed nominal Gram
values affect character coefficients only by signs.

Summing the bounded number of terms and averaging the nominal inputs
gives $\E_{\rawlaw^2}Q_n\le C_3\delta+o(1)$. Combine this with
\eqref{eq:product-four-corner} and
\eqref{eq:product-four-corner-typical}. The same $Q_n(O)$ bounds every
choice of separate weights, so the exceptional set is independent of
them.
\end{proof}

\begin{proof}[Proof of \Cref{thm:weak-product}]
If $r=1$, quantize the values of $W_n$ in $[-B,B]$ with mesh at most
$\eps/2$, choosing representatives in the same interval. The resulting
$S_n$ satisfies $\norm{W_n-S_n}_\infty\le\eps/2$ and is a linear
combination of at most $1+\lceil4B/\eps\rceil$ global level-set
indicators, with coefficients bounded by $B$. Both required errors
are at most $\eps/2$, since the reference measure is a probability
measure and the actual-query indicator and weights have bound one.
This includes $B=0$, when $S_n=0$. For $r\ge2$,
apply \Cref{lem:product-twisted-regularity} to a nontrivial split of
the full set of positions. Recursively apply it to group factors in
the structured terms until every factor is unary. Whenever a
remainder is produced, leave that whole term terminal; do not expand
its other factors. We obtain the exact decomposition, on the
unconditioned product reference,
\begin{equation}\label{eq:product-recursive-decomposition}
 W_n=S_n^{\mathrm{tw}}+\sum_{v\in\mathcal V_n}R_{v,n}.
\end{equation}
Each structured term is a product of unary global functions times
available Gram twists. Each terminal term has the form
\begin{equation}\label{eq:product-terminal-form}
 R_{v,n}(k)=c_{v,n}D_{v,n}(k_I,k_J)
       \prod_{H}f_{v,H,n}(k_H)\,\tau_{v,n}(k),
\end{equation}
where $I,J$ partition one group, the groups $H$ are disjoint outside
$I\cup J$, all outside functions are bounded, $\tau_{v,n}$ is a
product of available Gram signs, and $D_{v,n}$ has a specified small
twisted cut norm on $I|J$. The sup bounds, numbers of terms, and
partition complexities are all bounded independently of $n$.

Here is an explicit way to organize the tolerance choices. Along
any structured branch a split replaces one group by two, so at most
$r-1$ splits are needed before all groups are unary. At any fixed
depth the number of pending terms and all their coefficient and
factor bounds have already been bounded using previous choices.
Allocate at most $\eps/(2r)$ of each desired final error to that
depth, dividing this budget among its pending terms and their known
outside bounds. Choose the next cut tolerances small enough for the
reference estimates below and for the fourth roots in
\Cref{lem:product-terminal-remainder}. These choices determine a
finite bound on the next level's number of structured terms.
Continue for at most $r-1$ levels. The empirical tolerance $\varepsilon_{\mathrm{emp}}$ for
each terminal test is chosen at the same stage, sufficiently small
for its allotted fourth-root error. All choices are constants; none
requires a change in the construction parameters or in $M_0$.

On $\cK_n$, every Gram twist in $S_n^{\mathrm{tw}}$ equals one.
Removing them gives the function $S_n$ in
\eqref{eq:weak-product-approximant}. Its unary factors are indicators
of the finite partitions used in the recursion; coefficients and
term counts have the asserted bounds.

To prove the reference assertion, multiply a terminal term
\eqref{eq:product-terminal-form} by a product of arbitrary bounded
unary tests and expand the off-position Gram event defining $\nu_n$.
Fix all positions outside its remainder. Their group factors are
bounded constants. Gram phases internal to $I$ or $J$ and phases
involving a fixed outside position are separate bounded group
weights. The only other phases are available crossing twists.
Thus the integral is bounded by the twisted cut norm of $D_{v,n}$
times a constant depending only on its known bounds and the fixed
Gram normalization. The tolerance allocation makes the sum of these
errors at most $\eps$, proving \eqref{eq:weak-product-reference}.
This estimate is uniform over all unary reference test functions.

For the actual-query assertion, expand $\mathcal A_n$ into its
Gram and nominal $L,R$ characters. The number of terms is bounded.
Consider a terminal term and fix the outside position parameters.
All phases internal to $I$ or $J$, or involving an outside position,
can be absorbed into separate group weights. The remaining crossing
phases are a product $\psi_n(I,J)$ of precisely the form in
\Cref{lem:product-terminal-remainder}. It is independent of the
outside parameter values: a binary bit involving an outside
position has already been put into a separate weight. All original
unary weights, even when they depend on the orientation, can also be
included in these two weights.

Apply \Cref{lem:product-terminal-remainder}, scaling the known
outside bounds as necessary. Its common four-copy test is
independent of the fixed outside parameters, so one exceptional
orientation set controls every such fixation. Integrate the outside
parameters, sum the bounded binary-character expansion, and sum the
terminal terms. The chosen cut and empirical tolerances make the
total at most $\eps$ for sufficiently large $n$. The union of the
finitely many exceptional sets is exponentially small. Since each
four-copy estimate is uniform over its separate weights, the final
exceptional set is independent of all unary input weights and of
orientation-only filters. This proves
\eqref{eq:weak-product-actual}.
\end{proof}

\begin{corollary}[Factorized integration of the approximant]
\label{cor:weak-product-factorized-integration}
For the structured sum in \eqref{eq:weak-product-approximant} and
arbitrary bounded unary input weights, outside exponentially small
mixed mass,
\begin{align}\label{eq:weak-product-factorized-integration}
 &\E_z\mathcal A_n(O,z)S_n(K(O,z))\prod_jw_j(O,z_j)\notag\\
 &\qquad=
 \kappa\sum_{a=1}^{L}c_{a,n}
       \prod_j\E_{z_j}
        \bigl[f_{a,j,n}(K_j(O,z_j))w_j(O,z_j)\bigr]+o(1).
\end{align}
The error is uniform over the unary weights. An orientation-only
filter may be included before averaging this identity over leaves.
In particular the same statement applies to marked mixtures and
retained type/colour pairs as in \Cref{thm:weak-product}.
\end{corollary}

\begin{proof}
Apply \Cref{lem:binary-mixing} to each structured term, with
$f_{a,j,n}(K_j)w_j$ as its unary weights. The number of terms and
their bounds are fixed at the chosen approximation accuracy, so the
sum of the exponentially small errors tends to zero. The exceptional
set in that lemma is independent of these weights.
\end{proof}

\begin{remark}[What the two comparisons do not assert]
\label{rem:weak-product-filter-class}
The reference error in \eqref{eq:weak-product-reference} is weak
testing error, not $L^1$ error. The actual comparison allows arbitrary
orientation-dependent \emph{unary} parameter weights, and the
specified binary Gram/$L,R$ conditions. It does not allow an
unprescribed joint input filter to be inserted without proof.
For instance, two complementary bounded tuple densities determined
by a high-rank global bilinear bit can have nearly uniform unary
marginals and zero overlap. A two-element histogram net would not
exclude that example. The additional actual-query conclusion of
\Cref{thm:weak-product}, proved by the four-copy expansion, is what
rules it out for the successful-query format used here. All factors
added to the common signature lists by this theorem are global key
functions. Nominal input status bits remain weights, not common key
cells in \Cref{lem:unary-positivity}.
\end{remark}
